\section{From a coefficient field to a finite-cutoff reference observable}
\label{sec:reference-match}
The signed comparison must control the original finite-cutoff
observable, not only its coefficient vector. This requires a
simultaneous marked argument. We explain the finite-dimensional
matching and the exact replicated difference used in that argument.

Fix a degree $d$ and scales $\operatorname{ref}(d)\le m\le k<N$.
All paths below use the same finest scale $N$ and the same history;
$m$ is the opening scale and $k$ the matching scale. Let $Q$ represent a continuous finite-cutoff observable
$f$ at scale $k$ through the full marked local representation.
Suppose its coefficient-plus-remainder bound is at most $c>0$,
and put
\[
 E\ge4(1+C_k)c,\qquad
 G=e^{(s_k-s_m)/8},\qquad w>0.
\]
Assume the signed transition hypotheses of
Section~\ref{sec:reference-difference}, the one-scale marked
transport hypotheses of Section~\ref{sec:local-transfer}, and
the necessary cell separation and chamber locality at each scale.
The common affordability condition is
\begin{equation}\label{eq:reference-match-afford}
       32(1+C_k)w a_m^{-d}E/G\le\theta_m/2.
\end{equation}
Here $\theta_i$ are the single-source allowances; their growth
condition is the one used in \eqref{eq:reference-affordable}.
The coefficient paths must also satisfy the corresponding
one-step singleton cost conditions there and in the marked
transport lemma. No availability of these analytic hypotheses
is inferred from \eqref{eq:reference-match-afford} alone.

\subsection{An approximate coordinate basis}
Apply the signed reference construction first to $Q$. It gives a
label $b$ with $\mu(I_kb)\le2E$ and a path representing
$Q-\mathcal P_k(b)$. Its opening price is $2w$; the required
initial affordability follows from
\eqref{eq:reference-match-afford} because
$2\le32(1+C_k)$.

For each coefficient coordinate $j$, apply the same construction
to the primary field with vector $e_j$ and zero remainder.
Use opening price $8Ew$ and coefficient allowance $4(1+C_k)$.
Its affordability expression is exactly the left side of
\eqref{eq:reference-match-afford}. Obtain labels $b_j$.
The cone estimate at $k$, now with terminal bound $1$, gives
\[
 32(1+C_k)^2\mu(e_j-I_kb_j)\le2(1+C_k),
 \qquad \mu(e_j-I_kb_j)\le\tfrac14.
\]
Define $A z=\sum_jz_jI_kb_j$. Since $\mu$ is the coefficient
$\ell^1$ norm,
\[
 \mu(z-Az)\le\sum_j|z_j|\mu(e_j-I_kb_j)
                         \le\tfrac14\mu(z).
\]
It follows that $\mu(z)\le2\mu(Az)$. Thus $A$ is injective
and, in this finite-dimensional space, surjective. Solve
$As=I_kb$. Injectivity of $I_k$ gives
\begin{equation}\label{eq:reference-basis-expansion}
                   b=\sum_js_jb_j,
             \qquad\sum_j|s_j|\le2\mu(I_kb)\le4E.
\end{equation}
No uniform condition number of an arbitrarily chosen reference
basis is being assumed; the bound comes from the explicitly
constructed approximate coordinate vectors.

Let $\ell_j$ be the physically normalized primary field for
$e_j$ at scale $k$, with its local mask, and set
\[
                   \ell=-\sum_js_j\ell_j,
                         r=\mathcal B_{k,N}\ell,
\]
where $\mathcal B_{k,N}$ is the backward observable transfer to
the finest lattice. All $\ell_j$ lie in the same fixed stencil
window. Hence $\ell$ does too, and the room condition places
$r$ inside the centered chamber. This common window is essential
for the replicated transfer identity below.

\subsection{The full replicated difference}
We spell out that identity, keeping every mixed term. Write
$\mathcal T_{N,k}$ for the forward density transfer, adjoint to
$\mathcal B_{k,N}$ under the Haar priors. For copied fine
observables $O_i$ and copied coarse observables $F_i$, the
coefficient of a mark set $S$ after introducing the coarse
sources is
\begin{equation}\label{eq:reference-join}
 \sum_{A\subseteq S}
       \mathcal T_{N,k}\left(\rho_N
                    \prod_{i\in S\setminus A}(-O_i)\right)
                                      \prod_{i\in A}F_i.
\end{equation}
Here $\rho_N$ is the unnormalized Wilson density. The same
forward transfer is used for the whole fine product in each
summand. It has not been replaced by a product of transferred
factors.

For coarse functions in the common local window, the room
condition ensures that backward transfer respects their copied
product: $\mathcal B_{k,N}(\prod_{i\in A}F_i)$ is the product
of the corresponding copies of $\mathcal B_{k,N}F$.
Using adjointness in each summand of
\eqref{eq:reference-join}, its integral is therefore
\[
       \int\rho_N\prod_{i\in S}
                     (\mathcal B_{k,N}F_i-O_i)\,dU.
\]
This follows by the finite subset expansion of the product.
The interchange is justified by bounded local observables,
integrable fine densities and the transition-slice bounds.
Further forward transfers leave this integral unchanged.
Thus the coefficient of the complete mark set is the Wilson
partition function times the replicated difference moment.

The singleton bounds used here follow from the priced coefficient
basis and the remainder bounds. Explicitly, require the masked
basis evaluation bound with column constant $C_m\ge1$ and
remainder divisor at most one. If a coefficient field has vector
$q$ and remainder cost $t$, the triangle inequality gives
$|g|\le(1+C_m)\mu(q)+t$ before physical normalization.
The sum of the norms of its flagged profiles is at most $t$: sum
the absolutely summable forced majorants and use that every support
weight is at least one. The same argument bounds the ordinary
remainder by $t$; the primary term is bounded by
$(1+C_m)\mu(q)$ using the priced basis. If
$v a_m^{-d}(t+\mu(q))\le\theta_m$ and
$C_m\theta_m\le1$, these bounds give normalized singleton
allowances $2$ and $1$. Translations and conjugation preserve
them. This is the sharper bound for the matched coefficient paths;
it is separate from the general centered-field estimate.

Suppose the marked transport has reached scale $m$ with common
mark weight $v>0$, higher-array price one, and singleton bounds
\[
      v|g_i(U)|\le2,\qquad
          v\sum_X|q_i(X,U)|\le1.
\]
The integrated marked coefficient estimate
\eqref{ex-centered-sail}, with singleton allowances
$2$ and $1$, then bounds the complete moment by
\[
      v^M|\operatorname{moment}|\le
       \exp\bigl(|\Lambda_m|\mathcal V_m+5M\bigr).
\]
There are $M$ copied cells. Write
$S_0=\exp(|\Lambda_m|\mathcal V_m/M+5)$, the cell price.
Taking the $M$th root gives the centered root bound $S_0/v$.
This is an estimate for the complete replicated difference;
no estimate on its individual mixed summands is substituted.

\subsection{Combining the comparison paths}
The one-scale transport is stable under affordable finite linear
combinations of its single-source paths. Explicitly, if path
$j$ has positive opening price $v_j$ and coefficient $a_j$,
the condition
\begin{equation}\label{eq:reference-portfolio}
                           \sum_j w|a_j|/v_j\le1
\end{equation}
preserves every common normalized cost bound. All paths use the
same multiplicative weight evolution: at scale $i$ the opening
prices $v_j$ and $w$ are multiplied by the same positive factor.
Thus the fractions in \eqref{eq:reference-portfolio} are unchanged
at each intermediate scale. To verify the cost bound,
if $v_j\sum_X|F_{j,X}|\le C$, then
\[
 w\sum_X\left|\sum_ja_jF_{j,X}\right|
       \le\sum_j\frac{w|a_j|}{v_j}
                          v_j\sum_X|F_{j,X}|\le C.
\]
The same calculation applies to singleton values, elementary
transition costs and response majorants. Sum the dominated
indexed arrays using their summed majorants. The transition
identities are linear, so remain exact. Applying the marked
transport successively proves the claim for the whole interval.

Use the base difference path with coefficient $1$ and price
$2w$, and the coordinate paths with coefficients $-s_j$ and
price $8Ew$. Equation~\eqref{eq:reference-basis-expansion} gives
\[
                 \frac{w}{2w}
                     +\sum_j\frac{w|s_j|}{8Ew}\le1.
\]
At scale $k$, the reference terms cancel separately in the
ordinary and flagged profiles, because
$-b+\sum_j(-s_j)(-b_j)=0$.
For a conjugating copy, use the conjugate of every coefficient
$a_j$ in the portfolio; for an ordinary copy use $a_j$ itself.
Their absolute values coincide, so the cost proof is unchanged.
The cancellation in a conjugating copy is the conjugate of the
displayed identity. Thus cancellation holds for every copied profile.
The resulting ordinary field is the one for $Q$ plus $\ell$;
the flagged profile is exactly the one for $Q$. Introducing
$\ell$ into the marked representation of $f$ is multiplication
by its squarefree source exponential. The identity
\eqref{eq:reference-join}, with $O=-f$ and $F=\ell$, therefore
identifies the complete coefficient with the replicated moment
of $f+r$. The transport and complete-moment estimate give
\[
                           \mathcal R(f+r)\le S_0/w,
\]
where $\mathcal R$ is the centered root seminorm.

For each coordinate path the same argument, starting from the
terminal reference field and inserting $\ell_j$ at scale $k$,
gives
\[
 \mathcal R\bigl(\mathcal B_{k,N}\ell_j-\mathcal O(b_j)\bigr)
                                         \le S_0/(8Ew).
\]
Here $\mathcal O(b_j)$ denotes the fine observable attached to
that bare label by the calibrated terminal identification.
This application requires the terminal reference representation
and its singleton profile identity on the same history.
They remain among the analytic inputs to the matched construction.
The chamber seminorm inequality, backward-transfer linearity
and \eqref{eq:reference-basis-expansion} imply
\[
           \mathcal R\bigl(-r-\mathcal O(b)\bigr)
                  \le4E\frac{S_0}{8Ew}=\frac{S_0}{2w}.
\]
Finally $f-\mathcal O(b)=(f+r)+(-r-\mathcal O(b))$ lies in
the common chamber class, so its centered root is at most
$2S_0/w$. The matched reference observable has quota at most
$2E$, as proved above. These are exactly the two bounds used
in the cutoff matching argument.
